\section{Results}
\label{sec:results}

\subsection{Pilot training outcome (engineering success)}
The GPU pilot completed 1,641 optimizer steps across three full epochs. Epsilon
MSE declined monotonically:

\begin{table}[htbp]
\centering
\caption{Epoch-average epsilon MSE (GPU pilot, 27,169 parameters)}
\begin{tabular}{c c c}
\toprule
Epoch & Steps & Mean $\epsilon$ MSE \\
\midrule
1 & 547 & 0.219863 \\
2 & 547 & 0.180404 \\
3 & 547 & 0.172029 \\
\bottomrule
\end{tabular}
\label{tab:epoch-loss}
\end{table}

First-step MSE was 1.100, final-step MSE was 0.150. All losses and gradient norms
were finite. Training time was approximately 41 seconds (plus data loading).
Peak GPU allocation was 52.9~MiB (of 24~GiB). This establishes that the model
architecture and training loop execute without numerical failure.

\subsection{Sample quality (negative result)}
Despite declining training loss, \emph{generated sample magnitudes are physically
implausible}. The GPU pilot produced 8 finite 24-hour residual sequences spanning
$-16.91$ to $+7.94$ metres with SD 4.47~m. Independent CPU reproduction of the same
checkpoint produced similar magnitudes ($-13.08$ to $+16.28$~m, SD 4.92~m).

\begin{table}[htbp]
\centering
\caption{Pilot sample magnitudes vs.~training reference}
\begin{tabular}{l c c c}
\toprule
Quantity & Training & GPU samples & CPU reproduction \\
\midrule
Mean (m) & 0.108 & --- & --- \\
SD (m) & \textbf{0.215} & 4.474 & 4.922 \\
Min (m) & --- & $-16.910$ & $-13.079$ \\
Max (m) & --- & $+7.940$ & $+16.277$ \\
SD ratio vs.~train & 1.00 & \textbf{20.8} & \textbf{22.9} \\
\bottomrule
\end{tabular}
\label{tab:sample-magnitude}
\end{table}

The SD ratio of $\approx 22.95$ means that generated samples are roughly
\emph{23 times more spread out} than the training target distribution. This is
catastrophic scale explosion. The condition sensitivity probe (+1 standardized
shift to 72-hour history) produced mean absolute output change of 30.8
standardized units, confirming conditioning is active but not informative about
realistic magnitudes. The eight nearest-training-window RMSE probes ranged from
0.170 to 41.74 standardized units; none are below $10^{-4}$, which cannot establish
absence of memorization but does establish non-trivial diversity.

\subsection{Terminal-step amplification diagnosis}
The cosine schedule's final step ($t=63$) has $\beta_{63} \approx 0.999$ and
$\bar{\alpha}_{63} \approx 5.93 \times 10^{-7}$. This implies that the clean
reconstruction is amplified relative to the epsilon prediction by:
\begin{equation}
    \frac{1}{\sqrt{\bar{\alpha}_{63}}} \approx 1298.8.
    \label{eq:amplify}
\end{equation}
On the first eight fixed training windows with seed 20261028, terminal epsilon MSE
was approximately 0.005746 (epsilon $\sigma \approx \sqrt{0.005746} \approx 0.076$).
The clean reconstruction error scales as $\sigma_{\text{clean}} \approx
\sigma_\epsilon / \sqrt{\bar{\alpha}_{63}}$ (equation~\ref{eq:amplify}), giving
$0.076 \times 1298.8 \approx 99$~metres of implied clean RMSE. The observed implied
clean $\sigma$ at step 63 was 18.6~m (Table~\ref{tab:reverse-trace}), somewhat
smaller than this worst-case bound because the terminal-step amplification
applies only to the latest step's epsilon error accumulation, not an independent
worst case on every step. Regardless, the implication is unambiguous:
\emph{terminal alpha bar of order $10^{-7}$ amplifies any non-negligible epsilon
prediction error to catastrophic scale.}

\subsection{Oracle schedule comparison}
To isolate whether terminal schedule shape is a sufficient cause of the
magnitude failure, we run the \emph{same trained checkpoint} (27,169 parameters,
three epochs on cosine 64-step schedule) with four different inference
schedules without retraining. This is an oracle probe --- the checkpoint's
epsilon predictions are schedule-specific, so this is not a scientific fix
demonstration --- but it directly isolates whether terminal $\bar{\alpha}$ drives
sample magnitude.

\begin{table}[htbp]
\centering
\caption{Oracle schedule comparison --- same checkpoint, different inference}
\small
\begin{tabular}{l c c c c}
\toprule
Schedule & $\bar{\alpha}_T$ & Sample $\sigma$ (m) & Range (m) & Ratio \\
\midrule
Cosine 64 (historical) & $5.93\times10^{-7}$ & 4.922 & [$-$13.1, $+$16.3] & 22.95$\times$ \\
Linear 64 & 0.523 & 0.099 & [$-$0.16, $+$0.37] & 0.46$\times$ \\
Sigmoid 64 & 0.525 & 0.100 & [$-$0.17, $+$0.37] & 0.47$\times$ \\
\textbf{Cosine 1000$\rightarrow$64} & $\mathbf{5.89\times10^{-4}}$ & \textbf{0.297} & \textbf{[$-$0.69, $+$0.91]} & \textbf{1.38$\times$} \\
\midrule
\emph{Train residual} & --- & 0.215 & [$-$1.25, $+$1.53] & 1.00$\times$ \\
\bottomrule
\end{tabular}
\noindent\footnotesize Ratio: sample $\sigma$ / train $\sigma$. Amplification = $\sqrt{(1-\bar{\alpha}_T)/\bar{\alpha}_T}$.
\label{tab:oracle-schedules}
\end{table}

Linear and sigmoid schedules both produce physically plausible \emph{ranges} but
over-shrink the distribution ($\sigma = 0.10$~m, less than half the training standard
deviation). This is expected: the checkpoint was trained to predict epsilon under
the cosine schedule's noise levels, and those epsilon predictions do not transfer
to linear/sigmoid inference without retraining. Strikingly, \textbf{cosine 1000
subsampled to 64 steps} produces a physically reasonable sample: $\sigma = 0.297$~m
(1.38$\times$ train, well within the $0.1$--$5.0$ engineering scale screen),
range [$-$0.69, $+$0.91]~m covering most of the training distribution. The key
difference is terminal $\bar{\alpha}_{63} \approx 5.89\times10^{-4}$ (41.2$\times$
amplification) rather than $5.93\times10^{-7}$ (1299$\times$). This probe establishes
that \textbf{terminal schedule shape is a principal driver of the explosion
effect}: replacing a cosine schedule with $\bar{\alpha}_T \approx 10^{-7}$ by
one with $\bar{\alpha}_T \approx 10^{-4}$ reduces sample standard deviation by
an order of magnitude, without any retraining.

\subsection{Reverse-step trace}
Tracking per-step implied clean standard deviation during reverse sampling
reveals the explosion dynamics:

\begin{table}[htbp]
\centering
\caption{Per-step implied clean $\sigma$ during cosine reverse sampler}
\begin{tabular}{c c c c c}
\toprule
Step & $\bar{\alpha}_t$ & Amp & $\epsilon$ RMS & Implied $\sigma$ (m) \\
\midrule
63 & $5.93\times10^{-7}$ & 1299$\times$ & 1.140 & \textbf{18.64} \\
60 & $5.33\times10^{-3}$ & 13.7$\times$ & 1.278 & 5.76 \\
55 & $3.75\times10^{-2}$ & 5.1$\times$ & 1.082 & 5.05 \\
48 & $1.28\times10^{-1}$ & 2.6$\times$ & 0.863 & 5.09 \\
32 & $4.70\times10^{-1}$ & 1.1$\times$ & 0.702 & 4.98 \\
16 & $8.29\times10^{-1}$ & 0.5$\times$ & 0.582 & 4.94 \\
\bottomrule
\end{tabular}
\label{tab:reverse-trace}
\end{table}

The implied clean $\sigma$ remains stable at $\approx 5.0$~m for steps 0--48 (amplification
$\leq 2.6\times$), then \emph{explodes} at steps 55--63 as amplification dominates the
epsilon-prediction error that grows slowly throughout the reverse chain. This
pattern rules out a model-conditioning failure or a reversed-schedule-order bug:
the core issue is terminal-step schedule geometry interacting with imperfect
epsilon predictions at limited training budget.

\subsection{Engineering validation guards}
The repaired codebase adds nonfinite rejection at every reverse step and a
train-reference scale screen. The historical GPU samples fail this screen
(ratio 22.95 $\gg$ 5.0 threshold), as do the CPU reproductions. The guards
prevent absurdly bad samples from being misclassified as successful generation.
These are validation infrastructure, not a model repair.

\subsection{Generative model benchmark: Flow Matching vs.\ Diffusion-TS (GPU full-budget)}
We benchmark two additional generative frameworks against the same 96$\rightarrow$72
contract across all five admitted regions with two seeds each (GPU RTX3090, 15
epochs, Adam lr=1e-3, batch size 256, $\approx$28--39K parameters per model).

\begin{itemize}
    \item \textbf{M4 — Flow Matching / CNF}~\citep{lipman2023flow}. A continuous
    normalising flow that learns a single velocity-regression loss, no beta
    schedule and no variance annealing. Straight-line interpolation $x_t = (1-t)x_0
    + tx_1$ with Euler ODE integration. The VelocityTCN architecture mirrors the
    DiffusionTS residual denoiser (dilated TCN, 27,937 parameters) for fair
    architectural control.

    \item \textbf{M2 — Diffusion-TS decomposition}. Decomposes water level into
    trend (24h causal MA) + seasonal (sinusoidal regression on M2/K1 tidal
    periods) + residual. Trend and seasonal are extrapolated deterministically;
    the DDPM with cosine1000-subsample schedule operates on the residual
    component only (ResidualDenoiser, 39,425 parameters).
\end{itemize}

Both models operate on the same train-only standardized data, per-sample residual
z-score normalization with history statistics, and forecast 72-hour non-tidal
residuals. Quality evaluation uses the full \texttt{quality.py} battery:
Wasserstein-1, PSD deviation, autocorrelation RMSE, scale ratio, peak KS-test,
tail-exceedance ratio, and duplicate detection. The gate policy uses
$\text{scale} \in [0.5, 2.0]$, $\text{range} \in [0.3, 3.0]$, and
$\text{peak-}p > 0.01$.

\begin{table}[htbp]
\centering
\caption{Generative model benchmark (5 regions, 2 seeds each, test-set generated residuals)}
\small
\begin{tabular}{lc cccc c}
\toprule
Model & Region & Scale ratio & Peak ratio & q95 Gen/Real & q99 Gen/Real & Gate \\
\midrule
\multirow{5}{*}{M4 Flow Matching}
    & ny\_bight & 1.078 & 1.336 & 1.148/1.156 & 1.521/1.470 & INCONCLUSIVE \\
    & thames & 1.032 & 1.232 & 5.321/5.384 & 6.136/6.081 & INCONCLUSIVE \\
    & gulf & 0.953 & 1.356 & 0.697/0.695 & 1.004/0.951 & INCONCLUSIVE \\
    & charleston & 1.085 & 1.384 & 0.964/0.974 & 1.195/1.129 & INCONCLUSIVE \\
    & australia & 1.014 & 1.191 & 1.607/1.587 & 1.806/1.768 & INCONCLUSIVE \\
\midrule
\multirow{5}{*}{M2 Diffusion-TS (DDPM)}
    & ny\_bight & 0.087 & 26.2 & 11.1/1.156 & 16.9/1.470 & REJECTED \\
    & thames & 0.104 & 16.5 & 24.5/5.384 & 36.1/6.081 & REJECTED \\
    & gulf & 0.105 & 18.6 & 3.54/0.695 & 6.08/0.951 & REJECTED \\
    & charleston & 0.103 & 29.7 & 8.56/0.974 & 13.6/1.129 & REJECTED \\
    & australia & 0.107 & 14.3 & 6.95/1.587 & 10.1/1.768 & REJECTED \\
\bottomrule
\end{tabular}
\label{tab:gen-benchmark}
\end{table}

\textbf{Flow Matching produces near-perfect scale matches} across all five
regions: scale ratios of 0.95--1.12 (within $\pm$12\% of real), q95 and q99
quantiles within 3\% of real values. The INCONCLUSIVE gate decisions stem
exclusively from peak KS-test $p$-values (peak-magnitude distribution mismatch,
expected from single-model generation) and near-duplicate ratio warnings on the
thames/australia regions; neither triggers a hard rejection. Scale is correct.

\textbf{Diffusion-TS (DDPM-based) is catastrophically rejected everywhere.}
Scale ratios of 0.08--0.13 mean generated residuals are 8--12$\times$
\emph{under-spread} on the bulk distribution, while peak ratios of 14--30$\times$
mean extreme spikes explode beyond physical range. One seed even collapses to
NaN training loss on the gulf region. This occurs \emph{even with the
cosine1000-subsample schedule fix} (terminal $\bar{\alpha}_T \approx 5.89\times
10^{-4}$, 41$\times$ theoretical amplification). The schedule fix reduces
amplification by two orders of magnitude but does not eliminate the DDPM's
reverse-time accumulation of epsilon-prediction error.

The failure mechanism is distinct from the historical pilot's 1299$\times$
amplification: with $\bar{\alpha}_T \approx 6\times10^{-4}$, the amplification
term per reverse step is $\sqrt{(1-\bar{\alpha}_t)/\bar{\alpha}_t} \approx 41$,
but this amplifies the \emph{residual} component of a trend/seasonal
decomposition whose scale is already small relative to the full water level.
The decomposition itself works correctly; the DDPM reverse sampler on the
residual sub-space does not. This generalises the E6/E6-CPU finding:
\textbf{DDPM reverse-time amplification is not fully mitigated by schedule
tuning alone, even when amplification is reduced by two orders of magnitude.}

Flow Matching (CNF) fundamentally avoids this failure mode by construction: the
single velocity-regression loss trains a field that transports samples from a
simple Gaussian prior to the data distribution, and inference integrates the
learned ODE forward in time. There is no reverse chain and no per-step
amplification of accumulated error. The architecture is otherwise matched
(dilated TCN, same parameter budget, same training data), so the difference is
attributable to the \emph{generative framework}, not model capacity.

\subsection{A0 Real-only forecasting baselines}
We evaluate standard forecasting architectures on the held-out test set under
the 96$\rightarrow$72 contract.

\begin{table}[htbp]
\centering
\caption{A0 forecasting baselines on held-out test set}
\small
\begin{tabular}{ll cccc}
\toprule
Region & Model & MAE & RMSE & q95 Tail MAE & Peak MAE \\
\midrule
\multirow{4}{*}{ny\_bight}
    & Persistence & 0.170 & 0.228 & 0.333 & 0.301 \\
    & Seasonal & 0.154 & 0.214 & 0.310 & 0.150 \\
    & GRU & 0.112 & 0.161 & 0.340 & 0.152 \\
    & TCN & \textbf{0.112} & \textbf{0.158} & \textbf{0.337} & \textbf{0.146} \\
\midrule
\multirow{4}{*}{thames}
    & Persistence & 0.641 & 0.798 & 1.027 & 0.959 \\
    & Seasonal & 0.442 & 0.579 & 0.653 & 0.239 \\
    & GRU & \textbf{0.207} & 0.286 & 0.592 & 0.389 \\
    & TCN & 0.218 & \textbf{0.297} & \textbf{0.576} & \textbf{0.355} \\
\midrule
\multirow{4}{*}{gulf}
    & Persistence & 0.137 & 0.189 & 0.248 & 0.205 \\
    & Seasonal & 0.131 & 0.186 & 0.253 & \textbf{0.131} \\
    & GRU & \textbf{0.098} & \textbf{0.142} & 0.264 & 0.122 \\
    & TCN & 0.103 & 0.150 & \textbf{0.280} & 0.126 \\
\midrule
\multirow{4}{*}{charleston}
    & Persistence & 0.176 & 0.222 & 0.221 & 0.308 \\
    & Seasonal & 0.149 & 0.191 & 0.182 & \textbf{0.098} \\
    & GRU & \textbf{0.095} & \textbf{0.129} & 0.178 & 0.107 \\
    & TCN & 0.096 & 0.130 & \textbf{0.185} & 0.101 \\
\midrule
\multirow{4}{*}{australia}
    & Persistence & 0.051 & 0.066 & 0.074 & 0.079 \\
    & Seasonal & 0.055 & 0.070 & 0.084 & 0.054 \\
    & GRU & 0.046 & \textbf{0.061} & 0.112 & 0.061 \\
    & TCN & \textbf{0.046} & 0.060 & \textbf{0.107} & \textbf{0.056} \\
\bottomrule
\end{tabular}
\label{tab:a0-baselines}
\end{table}

GRU and TCN reduce MAE by 30--50\% over persistence and seasonal baselines on
all regions except the low-variance australia region (where all methods converge
to $\approx$0.05~m MAE). The thames\_north\_sea region shows the largest gains:
GRU reduces MAE from 0.64 (persistence) to 0.21, a 68\% improvement. Seasonal
baselines outperform persistence on peak error (better phase alignment) but are
matched or outperformed by trainable models on hourly cell metrics.
